% Section 7 -- Conclusion
% Drafted 12 November 2026; edited by the authors and revised 13 November 2026.
\section{Conclusion}\label{sec:conclusion}

% 1. The question
In England, older owners live in less energy-efficient homes than younger owners. This paper asks why. While keeping the preferences of older and younger owners the same, it considers two channels through which age could affect investment in efficiency. The first is the horizon channel: an older owner has fewer years to collect the return on a renovation. The second is the purchase window: a buyer gets efficiency with the home for a small premium and recovers it when she sells, but this route is open only to owners who move.


% 2. The answer
By solving and estimating a lifecycle model on English data we find that the age gradient comes
from the purchase window, not from the horizon channel. At the estimated parameters, the horizon is
short at every age, because efficiency wears off long before most owners die. Owners of every age
choose alike when they buy, and few renovate in place at any age. However, owners buy less often as they age. When buying an efficient home is
ruled out, the gap in the efficient share between the youngest and the
oldest owners falls from about ten percentage points to about one. The
result holds even when efficiency lasts 25 years instead of six. The horizon
then falls with age, but purchases still account for most of the decline.

% 3. What it changes: ageing and policy
The age gradient is therefore a purchase gradient.
There are two implications associated with this finding. First, longer lives do little to raise investment
in efficiency. As the population ages, a growing share of homes is held
by owners who rarely buy, and these homes are seldom upgraded. Second, the route
through which a policy acts decides whom it reaches. For example, a subsidy for
renovation reaches older owners, because they are primarily the ones who rely on
renovating in place. This narrows the age gap. A standard imposed at the
point of sale raises efficiency the most, but more so among the younger households, and steepens the age gradient. Also, by
making a purchase dearer it slows the turnover through which efficiency
enters the stock. Reducing the transaction costs associated with buying and selling opens more purchase opportunities but
weakens the incentive to renovate in place. A higher price of energy raises
efficiency at every age in about the same proportion and leaves the
gradient in place, while a larger price premium on efficient homes closes
the age gap by making the homes of young owners less efficient, not those
of older owners more efficient.

% 4. Beyond homes
Although we study energy efficiency, the mechanism is not specific to energy. Whenever the quality of a durable
is cheapest to change at the moment it is bought, the quality a household
holds depends on how often it transacts. Hence, households that rarely
transact hold quality chosen long ago. Homes are rarely replaced but often
resold, and building codes reach them only when they are new, so for the
existing stock the sale is the main moment at which efficiency improves.
